% Short running form; the full title is too wide for the head's right box.
\runninghead{The King Bound Before Pilate}
\section{The King Bound Before Pilate, and the Lamb Upon the Throne}
\label{sec:opening}

The Mass of Our Lord Jesus Christ the King begins in heaven and comes down, at
its Gospel, to a prisoner. Its Introit takes up the song of the angels around
the throne in the fifth chapter of the Apocalypse: ``The Lamb that was slain is
worthy to receive power and divinity and wisdom and strength and honour,'' and
adds from the first chapter, ``To him be glory and empire for ever and ever.''
Its verse is the opening prayer of Psalm~71 (72), ``Give to the king thy
judgment, O God, and to the king's son thy justice.'' The Collect confesses
that God willed to restore all things in his beloved Son, \latin{universorum
Rege}, the King of all, and asks that the families of the nations, scattered
by the wound of sin, \latin{peccati vulnere disgregatae}, be brought under his
most gentle rule, \latin{eius suavissimo subdantur imperio}. The Epistle is
St~Paul's thanksgiving to the Father ``who hath delivered us from the power of
darkness and hath translated us into the kingdom of the Son of his love''
(Col 1:13), and his hymn to that Son: the image of the invisible God, the
firstborn of every creature, in whom all things were created, the head of the
body, the Church, the firstborn from the dead, through whom the Father
reconciles all things, ``making peace through the blood of his cross.'' The
Gradual returns to Psalm~71: ``he shall rule from sea to sea, and from the
river unto the ends of the earth \dots\ all kings of the earth shall adore him:
all nations shall serve him.'' The Alleluia sings the second half of Daniel's
vision of one like a son of man: ``his power is an everlasting power that shall
not be taken away: and his kingdom that shall not be destroyed.''

The Gospel is the first questioning of Jesus by Pilate inside the governor's
hall (Jn 18:33--37). ``Art thou the king of the Jews?'' Jesus answers with a
question, ``Sayest thou this thing of thyself, or have others told it thee of
me?'' Pilate: ``Am I a Jew? Thy own nation and the chief priests have
delivered thee up to me. What hast thou done?'' Then the words on which the
feast turns: ``My kingdom is not of this world. If my kingdom were of this
world, my servants would certainly strive that I should not be delivered to
the Jews: but now my kingdom is not from hence.'' ``Art thou a king then?''
``Thou sayest that I am a king. For this was I born, and for this came I into
the world; that I should give testimony to the truth. Every one that is of
the truth heareth my voice.'' The pericope stops there, before Pilate's ``What
is truth?'' The Offertory gives the Father's word to his Son in the second
psalm, ``Ask of me, and I will give thee the Gentiles for thy inheritance, and
the utmost parts of the earth for thy possession.'' The Secret offers
\latin{hostiam \dots\ humanae reconciliationis}, the victim of mankind's
reconciliation, and asks that the Son, immolated in the present sacrifice,
grant all nations the gifts of unity and peace. The proper Preface gives
thanks that the Father anointed his only-begotten Son with the oil of
gladness, eternal Priest and King of all things, so that, offering himself
on the altar of the cross, he might accomplish the mysteries of man's
redemption, and, all creatures being subjected to his rule, deliver to the
Father's majesty ``an eternal and universal kingdom, a kingdom of truth and
life; a kingdom of holiness and grace, a kingdom of justice, love and peace''
(in the English of the Lasance missal). The Communion is two half-verses of
Psalm~28 (29): ``the Lord shall sit king for ever \dots\ the Lord will bless
his people with peace.'' The Postcommunion asks that those who have received
the food of immortality, and who glory to fight under the standards of
Christ the King, may reign with him for ever in the heavenly seat.

The Missal prints this Mass in the Sanctorale, among the feasts of October,
under the heading \latin{Dominica ultima octobris}, \latin{D.~N. Iesu Christi
Regis}, of the first class (printed pp.~697--698). Pius~XI instituted the feast
by the encyclical \work{Quas primas} of 11~December 1925 and appointed it to the
last Sunday of October, the Sunday that immediately precedes All Saints
(n.~28); the Sacred Congregation of Rites approved its Mass and Office the next
day, and the Mass was first printed with that decree in the \work{Acta
Apostolicae Sedis}. On Sunday 25~October 2026, which is both the last Sunday of
October and the Sunday before All Saints, the feast takes the place of the
Twenty-second Sunday after Pentecost with all its rights, and no commemoration
of the Sunday is made (RG~17~d; RG~112~b), as the formulary itself directs
after each of its three orations. The commemoration of Saints Chrysanthus and
Daria, inscribed on this date, is omitted (RG~111~a). The colour is white; the
Gloria and the Credo are said, and the Preface is the feast's own (RGMR~492).

Three questions rise from these texts, and each governs one of the readings
that follow. The first is what the Lamb receives. The Introit says that the
slain Lamb is worthy \latin{accipere \dots\ divinitatem}, to receive divinity,
and the Offertory has the Father bid his Son ask for the nations; of the
eternal Son both seem to say that he receives what he could never lack. St
Hilary of Poitiers answers from the second psalm that the risen Son asks back
the glory he had before; St Thomas Aquinas distinguishes a dominion that is
his by birth from a dominion he merited by prayer and passion; St Bede sees in
the Lamb who dies and the Lion who conquers one Lord, King and Priest by his
self-offering. The second is where the kingdom is, and whom it subjects. The
bound prisoner says his kingdom is not of this world and not from hence, while
the Gradual sings that all kings shall adore him and the Collect asks that the
nations be subjected to him. St Augustine answers that the kingdom is not
\emph{of} this world but is \emph{in} it, gathered out of every nation; St John
Chrysostom, that its empire is from above and needs no servants to defend it;
St Robert Bellarmine, that the kings are bidden to subject themselves freely
now and will all bend at the judgment. The third is how the kingdom ends. The
Preface says the Son will hand the kingdom over to the Father, the Alleluia
that his power shall not be taken away, the Communion that he shall sit King
for ever. St Augustine answers that the kingdom handed over is the faithful,
brought to the contemplation of God; St Hilary, in his own words, that the Son
delivers to God not his power but us, made his kingdom by the glorifying of his
body. The three readings are complementary: they say how he reigns as man, what
kind of kingdom it is and how it claims the nations, and where it is going.

\begin{maptable}
Introit &
\latin{Dignus est Agnus, qui occisus est}; Apoc 5:12; 1:6; verse Ps 71:1 (72:1) &
The angels' acclamation of the slain Lamb as worthy to receive power and
divinity; glory and empire to him for ever; the prayer that God give his
judgment to the King and the King's Son.\\
Collect &
\latin{Omnipotens sempiterne Deus, qui in dilecto Filio tuo, universorum Rege} &
God willed to restore all things in his Son, King of all; that the families
of the nations, scattered by sin, be subjected to his most gentle rule.\\
Epistle &
Col 1:12--20 &
Thanks to the Father who has transferred us into the kingdom of his beloved
Son: the image of God, firstborn of creation and from the dead, head of the
Church, who makes peace through the blood of his cross.\\
Gradual &
\latin{Dominabitur a mari usque ad mare}; Ps 71:8, 11 (72:8, 11) &
Dominion from sea to sea and to the ends of the earth; all kings adore him,
all nations serve him.\\
Alleluia &
\latin{Potestas eius, potestas aeterna}; Dan 7:14 &
An everlasting power that shall not be taken away, a kingdom that shall not
be destroyed.\\
Gospel &
Jn 18:33--37 &
Bound before Pilate, Jesus says his kingdom is not of this world and not from
hence; he is a king, born to bear witness to the truth, heard by all who are
of the truth.\\
Offertory &
\latin{Postula a me}; Ps 2:8 &
The Father to the Son: ask, and I will give thee the nations for thy
inheritance and the ends of the earth for thy possession.\\
Secret &
\latin{Hostiam tibi, Domine, humanae reconciliationis offerimus} &
The victim of mankind's reconciliation is offered; that Christ, immolated in
this sacrifice, may give all nations unity and peace.\\
Preface &
\latin{Qui unigenitum Filium tuum \dots\ oleo exsultationis unxisti}; Missal
p.~291 &
The Son anointed eternal Priest and King of all, offering himself on the
cross, subjecting all creatures and handing over to the Father an eternal,
universal kingdom of truth and life, holiness and grace, justice, love and
peace.\\
Communion &
\latin{Sedebit Dominus Rex in aeternum}; Ps 28:10--11 (29:10--11) &
The Lord shall sit King for ever; the Lord will bless his people with peace.\\
Postcommunion &
\latin{Immortalitatis alimoniam consecuti} &
Fed with the food of immortality, those who glory to fight under the
standards of Christ the King ask to reign with him for ever in heaven.\\
\end{maptable}

The same page prints two further chants for other days: a Tract for votive
Masses of Christ the King said after Septuagesima, and an Alleluia of two
verses for Paschaltide. The feast itself, on the last Sunday of October, falls
in neither season, so neither is ever said on it. Both are given with their
rubrics at the end of the appointed texts.
